% !TITLE 送元二使安西
% !AUTHOR 王维
% !DYNASTY 唐
% !GENRE 七言绝句
% !ORDER 42

\begin{poem}
渭城朝雨浥轻尘\textsuperscript{1}，客舍青青柳色新\textsuperscript{2}。
劝君更尽一杯酒\textsuperscript{3}，西出阳关\textsuperscript{4}无故人。
\end{poem}

\begin{annotations}
\notetext{1}{渭城朝雨浥轻尘：渭城（今咸阳）的清晨小雨，湿润了地上的尘土。浥（yì），湿润。}
\notetext{2}{柳色新：柳树被雨洗过，青翠欲滴。“柳”谐音“留”，古人折柳送别。}
\notetext{3}{更尽一杯酒：再喝完这最后一杯。不是劝酒，是想多留对方一刻。}
\notetext{4}{阳关：通往西域的边关，在今甘肃敦煌西南。出了阳关便是茫茫大漠，再见不到熟人了。}
\end{annotations}

\begin{appreciation}
《送元二使安西》是中国文学史上最著名的送别诗之一。后来它被谱成曲子反复传唱，一段旋律叠三遍，就叫《阳关三叠》——送别要唱三遍才肯放下，跟这首诗的道理是一样的。

渭城朝雨浥轻尘——渭城就是今天的咸阳，清晨一场小雨，正好把地上的浮尘润湿了。浥是沾湿的意思。你可别小看这场雨：尘土压住了，路好走，这是适合上路的好天气。可越是好天气，越说明分别是真的要来了。客舍青青柳色新——旅店边的柳树被雨洗得发绿。"柳"和"留"同音，古人送行要折一枝柳给对方，意思是：别走太快，再留一留。

劝君更尽一杯酒，西出阳关无故人——请你再喝掉这一杯。阳关在敦煌附近，出了阳关就是茫茫大漠，再也遇不见一个熟人。所以这一杯不是劝酒，是挽留：能多坐一刻是一刻。王维没写眼泪，没写拥抱，全诗就这么一句大白话，却比哭天抢地都沉。千言万语，都在那杯酒里了。

还记得我们读过的《行行重行行》吗？两汉那位女子对游子说"弃捐勿复道，努力加餐饭"——说多了都是伤心，只盼你在外头好好吃饭。从"努力加餐饭"到"劝君更尽一杯酒"，隔了七八百年，中国人送别时的关心一点没变：我留不住你，只愿你路上好走。

历史上元二几乎什么都没留下，史书里翻不到他的名字。可因为这首诗，一千二百年来的读者都记住了他——被一杯酒记住，也算没有白来人间一趟。

你以后也会送朋友走，也会被人送走。哥哥没什么可嘱咐的，跟那位两汉的女子一样：出门在外，记得好好吃饭，到了地方，报个平安。
\end{appreciation}
